% Online Resource 1 of "Quantum Optimization for Adaptive Low-Rank CP Decomposition".
% Build after `python main.py tables`, from this folder:
%   pdflatex online_resource_1.tex && pdflatex online_resource_1.tex
\documentclass[10pt]{article}
\usepackage[margin=2cm]{geometry}
\usepackage{amsmath,amssymb}
\usepackage{booktabs,longtable,pdflscape}
\usepackage[T1]{fontenc}
\usepackage{lmodern}
\usepackage[hidelinks]{hyperref}
\setlength{\LTpre}{4pt}\setlength{\LTpost}{4pt}
\newcommand{\T}[1]{\input{../tables/revision/#1}}

\title{Online Resource 1\\[4pt]\large Quantum Optimization for Adaptive Low-Rank CP Decomposition\\
\normalsize Full result tables}
\author{Van Tien Nguyen and Panagiotis Markopoulos\\\small The University of Texas at San Antonio}
\date{}

\begin{document}
\maketitle

\section*{S1\quad Conventions}
Every table below is generated from the per-run records of the code repository
(\url{https://github.com/vantnprof/QALR-CPD}, folder \texttt{results/revision/}) by the command \texttt{python main.py tables}.
Entries are mean \(\pm\) standard deviation over seeds; \(n\) is the number of runs.

\begin{itemize}
\item \emph{Penalty.} \(\mu\) is dimensionless; every solver receives the QUBO with
\(\mu_{\mathrm{eff}}=\mu\|\mathcal X\|_F^2/R_{\max}\). Penalty grid of the selection protocol (Sec.~3.7): \(0.005\),
\(0.01\), \(0.02\), \(0.05\), \(0.1\), \(0.15\), \(0.2\), \(0.3\), \(0.4\), \(0.5\), \(0.65\), \(0.8\), \(1\), \(1.5\), \(2\).
A plateau is selected only if it spans at least three consecutive grid values; \(\mu\) is then its middle value (the
upper of the two middle values for an even number).
\item \emph{Datasets.} \texttt{synth\_default}: default synthetic benchmark (\(20^{4}\), \(R_{\mathrm{true}}=6\),
\(\rho=0.5\), \(\sigma=10^{-2}\)); \texttt{amino}: Amino Acids; \texttt{kinetic}: X60t (error over observed entries);
\texttt{covid}: COVID-19 systems serology; \texttt{snr}\(s\)\texttt{dB\_rho0.5}: exact-SNR grid;
\texttt{rho}\(r\)\texttt{\_snr}\(s\)\texttt{dB}: SNR-matched collinearity designs; \texttt{noise\_}\(\sigma\) and
\texttt{rho\_}\(r\): noise and collinearity sweeps of the first version.
\item \emph{Methods.} As in the manuscript: ALS at fixed rank; QALR-CPD with exact search (ES), bit flipping (BF),
multi-start bit flipping (MS-BF), simulated annealing (SA), tabu search (Tabu) or QAOA; in-loop magnitude pruning
(Mag-IL) and its \(z\)-form variant (Diag-x); post-hoc truncation without (TopK-PH) and with 50 refit sweeps
(TopK-PH+refit); in-loop truncation matched call by call to ES (TopK-IL).
\item \emph{Excess over \(\varepsilon^\star\)}: relative error minus the realized noise floor
\(\|\mathcal N\|_F/\|\mathcal X\|_F\), in percent of the floor. \(\hat R=R_{\mathrm{true}}\): share of seeds with exact
rank recovery.
\end{itemize}

\section*{S2\quad Controlled ablation (manuscript Table 9)}
\textbf{Table S1.} Default synthetic benchmark, all three penalties of the first version.
\T{tableR1_ablation_synthetic.tex}
\textbf{Table S2.} Real tensors, all three penalties of the first version.
\T{tableR1_ablation_real.tex}

\section*{S3\quad COVID-19 rows from \(R_{\max}=20\) (manuscript Table 10)}
\textbf{Table S3.}
\T{table7_covid_fixed_rows.tex}

\section*{S4\quad Exact-SNR grid and SNR-matched collinearity (manuscript Tables 7 and 8)}
\textbf{Table S4.} Every design at \(\mu\in\{0.1,0.3,0.5\}\).
\T{tableR2-7_snr_and_rho.tex}

\section*{S5\quad Penalty plateaus (manuscript Tables 14 and 15)}
\textbf{Table S5.} Widest stable plateau for every dataset, \(R_{\max}\) and selection rule. The number of grid values a
plateau spans is the number of grid values in its \(\mu\) range (Sec.~S1); plateaus of fewer than three values lead to
abstention.
\T{tableC1_plateaus.tex}

\section*{S6\quad Noise-calibrated starting penalty}
\textbf{Table S6.} \(\mu_0=\kappa\,\hat\sigma^2(\sum_n I_n)R_{\max}/\|\mathcal X\|_F^2\) and the rank exact search
selects at \(\mu_0\).
\T{tableC1_sigma_rule.tex}

\section*{S7\quad Core consistency}
\textbf{Table S7.} CORCONDIA at the plateau rank and its neighbours.
\T{tableC1_corcondia.tex}

\section*{S8\quad Solvers on identical saved QUBOs (manuscript Table 16 and Fig. 13)}
\begin{landscape}
\textbf{Table S8.} Main bank.
{\footnotesize\setlength{\tabcolsep}{4pt}\T{tableR2_solver_bank.tex}}
\end{landscape}
\textbf{Table S9.} Scaling bank: share of instances solved to optimality, by number of QUBO variables.
\T{tableR3-2_scaling.tex}

\section*{S9\quad QAOA under the uniform non-empty rule (manuscript Table 18)}
\textbf{Table S10.}
{\small\setlength{\tabcolsep}{4pt}\T{tableR2-6_qaoa_noise.tex}}

\section*{S10\quad First-version sweeps under one protocol (manuscript Tables 20 and 21)}
\textbf{Table S11.}
\T{tableR4-1_sweeps.tex}

\section*{S11\quad Split-half check of the selection protocol (manuscript Sec. 4.4)}
\textbf{Table S12.} Selection on one half of the seeds of each distinct synthetic design, exact recovery at the
selected \(\mu\) on the other half.
\T{tableS_heldout.tex}

\section*{S12\quad Hardware and noise-model repeats (manuscript Table 19)}
\textbf{Table S13.} One saved instance (Amino Acids, \(R=12\), \(\mu=0.1\), first pruning call), five repeats of 1024
shots per backend; angles optimized on the noiseless simulator. Noise models: fake backends of qiskit-ibm-runtime
0.50.0. Job identifiers are in \texttt{results/revision/hardware/hardware\_repeats.csv}.

\smallskip
{\footnotesize\setlength{\tabcolsep}{4pt}\T{tableS_hardware.tex}}

\section*{S13\quad Cost of one pruning call (manuscript Table 17)}
\textbf{Table S14.}
\T{tableR4-5_runtime.tex}

\end{document}
